\section{Declaration on the use of AI}
\note{
When considering the use of generative artificial intelligence (AI) tools for the preparation of the proposal, it is imperative to exercise caution and careful consideration. The AI-generated content should be thoroughly reviewed and validated by the applicants to ensure its appropriateness and accuracy, as well as its compliance with intellectual property regulations. Applicants are fully responsible for the content of the proposal (even those parts produced by the AI tool) and must be transparent in disclosing which AI tools were used and how they were utilized.

Specifically, applicants are required to:
\begin{compactitem}
    \item Verify the accuracy, validity, and appropriateness of the content and any citations generated by the AI tool and correct any errors or inconsistencies.
    \item Provide a list of sources used to generate content and citations, including those generated by the AI tool. Double-check citations to ensure they are accurate and properly referenced.
    \item Be conscious of the potential for plagiarism where the AI tool may have reproduced substantial text from other sources. Check the original sources to be sure you are not plagiarizing someone else's work.
    \item Acknowledge the limitations of the AI tool in the proposal preparation, including the potential for bias, errors, and gaps in knowledge.
\end{compactitem}
}

%%% Local Variables:
%%% mode: latex
%%% TeX-master: "master"
%%% TeX-PDF-mode: t
%%% End:
